\documentclass[11pt,a4paper]{article}
\usepackage[utf8]{inputenc}
\usepackage[T1]{fontenc}
\usepackage[ngerman]{babel}
\usepackage[margin=1in]{geometry}
\usepackage{xcolor}
\usepackage{hyperref}
\definecolor{primarycolor}{RGB}{33,37,41}
\definecolor{secondarycolor}{RGB}{0,102,204}
\hypersetup{colorlinks=true,linkcolor=secondarycolor,urlcolor=secondarycolor,citecolor=secondarycolor}
\setlength{\parindent}{0pt}
\setlength{\parskip}{0.85em}
\pagestyle{empty}
\begin{document}
{\bfseries\color{primarycolor} Katarzyna Wiśniewska}\\
Habsburgerstraße 41, 79104 Freiburg im Breisgau\\
k.wisniewska@example.com \textbar{} +49 761 555 0412

Sozialstation Dreisamtal gGmbH

\textbf{\color{primarycolor} Bewerbung als Pflegehelferin in Teilzeit mit 75 \%}

Sehr geehrte Damen und Herren,

seit April 2020 arbeite ich als Pflegehelferin in der stationären Altenpflege. Für Ihren ambulanten Pflegedienst bringe ich fundierte Erfahrung in der Grundpflege, die Qualifikation Behandlungspflege LG1 und einen ruhigen, geduldigen Umgang mit Menschen mit Demenz mit. Ihr Leitsatz „Pflege mit Zeit und Respekt“ passt zu meiner Arbeitsweise: zuverlässig helfen und auf den einzelnen Menschen eingehen.

Im Seniorenzentrum Am Lindenhof unterstütze ich pro Schicht 10 bis 12 Bewohnerinnen und Bewohner bei Körperpflege, Mobilisation, Lagerung und Essen. Zur Behandlungspflege nach LG1 gehören Blutdruck- und Blutzuckermessung, Kompressionsstrümpfe und das Reichen von Medikamenten nach Plan. Ich dokumentiere sicher in Vivendi PD und verfüge über ein Deutschzertifikat auf Niveau B2. Bei einem Projekt zur Sturzprophylaxe habe ich mit der Wohnbereichsleitung Sturzprotokolle ausgewertet und Maßnahmen umgesetzt; im folgenden Halbjahr gingen die Stürze im Wohnbereich von 14 auf 8 zurück. Zusätzlich leite ich wöchentlich eine Singgruppe für Menschen mit Demenz.

Meine bisherige Pflegeerfahrung stammt aus dem stationären Bereich; mit MediFox habe ich noch nicht gearbeitet. Ihre zweiwöchige begleitete Einarbeitung bietet einen passenden Rahmen für den Einstieg in die ambulante Pflege. Vor meiner Tätigkeit in Deutschland habe ich als Krankenschwester in Polen gearbeitet. Mein Anerkennungsverfahren als Pflegefachfrau läuft; derzeit bewerbe ich mich ausdrücklich als Pflegehelferin. Die Kenntnisprüfung ist für März 2027 vorgesehen.

Ich wünsche mir einen Stellenumfang von 75 \% und kann an jedem zweiten Wochenende arbeiten. Den geforderten Führerschein Klasse B besitze ich derzeit noch nicht: Die Theorieprüfung ist bestanden, die praktische Prüfung ist für Januar 2027 vorgesehen. Gerne möchte ich mit Ihnen klären, ob unter diesen Voraussetzungen eine Beschäftigung möglich ist. Über eine Einladung zum persönlichen Gespräch freue ich mich.

Mit freundlichen Grüßen\\[0.5em]
Katarzyna Wiśniewska
\end{document}